\chapter{Security Architecture and Hardening}

\section{Introduction}
This chapter presents the security architecture integrated into \projectname{}. The goal is not to claim a complete enterprise identity-management or zero-trust program, but to provide an implemented hardening baseline for service identity, browser access, exposed interfaces, remote operations, credential separation, and persisted runtime data. The controls are grounded in established secure-service mechanisms \cite{mit_kerberos_docs,apache_hadoop_docs,openssh_docs,cryptsetup_docs}.

\section{Security Objectives}
The security work of the project follows six main objectives:
\begin{itemize}
    \item authenticate critical distributed services rather than relying on implicit trust;
    \item secure browser-facing administrative or analytical interfaces;
    \item protect exposed endpoints through encryption in transit;
    \item reduce the operational risk of remote orchestration actions;
    \item separate service, administration, and demonstration identities and credentials;
    \item strengthen the protection of service data persisted on the virtual machines.
\end{itemize}

\begin{table}[!htbp]
    \centering
    \footnotesize
    \renewcommand{\arraystretch}{1.16}
    \arrayrulecolor{tablelineblue}
    \rowcolors{2}{tablebodyblue}{white}
    \begin{tabularx}{\textwidth}{|Y|Y|Y|}
        \hline
        \rowcolor{tableheaderblue}
        \textcolor{white}{\textbf{Security objective}} & \textcolor{white}{\textbf{Implemented mechanism}} & \textcolor{white}{\textbf{Protected scope}} \\
        \hline
        Distributed identity control & Kerberos-secured Hadoop runtime & Hadoop services and related secure interactions \\
        \hline
        Protected web access & SPNEGO + HTTPS for HDFS browser endpoints & NameNode and DataNode administrative web paths \\
        \hline
        Secured analytical exposure & authenticated Trino gateway over TLS & browser-based SQL access path \\
        \hline
        Protected orchestration path & SSH-based remote Spark submission & control-plane to VM2 operational execution \\
        \hline
        Operational identity separation & distinct service, administration, and demonstration identities with environment-scoped credentials & non-interactive services and human-operated access paths \\
        \hline
        Encrypted persisted runtime data & optional LUKS-backed runtime storage support & service data at rest on the Linux nodes when enabled \\
        \hline
    \end{tabularx}
    \caption{Main security objectives and corresponding implementation controls}
    \label{tab:chapter7-security-objectives}
\end{table}
\FloatBarrier

\section{Kerberos-Secured Hadoop Services}
Kerberos is the core identity mechanism used to secure the Hadoop side of the platform. In a distributed environment, services such as the NameNode and DataNodes should not trust requests only because they originate from another cluster component. They must validate the identity of the caller explicitly \cite{mit_kerberos_docs,apache_hadoop_docs}.

In the implementation, the Kerberos KDC and administration service are hosted on VM2, which is already the coordination-heavy node of the data plane. This is consistent with the overall architecture: the node that centralizes namespace and metadata authority also centralizes identity management for the secure Hadoop runtime.

The practical contribution of Kerberos in this project is threefold:
\begin{itemize}
    \item service principals and keytabs provide authenticated identities for non-interactive Hadoop services;
    \item it allows secure service-to-service interactions in the distributed cluster;
    \item it provides the foundation for browser-based negotiated authentication through SPNEGO.
\end{itemize}

The Hadoop runtime therefore uses explicit authentication rather than relying only on open internal trust between cluster components.

\begin{figure}[H]
    \centering
    \begin{tikzpicture}[x=1cm,y=1cm,font=\small]
  \node[ado box={6.60cm}{adotealfill}, minimum height=1.25cm, inner sep=5pt] (kdc) at (0,0)
    {\textbf{\normalsize Kerberos KDC --- VM2}\\[-1pt]
     \footnotesize principal administration and service-ticket issuing};

  \node[ado box={8.20cm}{adograyfill}, minimum height=1.25cm, inner sep=5pt] (requestor) at (0,-2.05)
    {\textbf{\normalsize Authenticated requestors}\\[-1pt]
     \footnotesize Spark/Hive clients and peer Hadoop services};

  \draw[ado arrow, draw=black!60, line width=1.05pt]
    ($(requestor.north)+(-1.20,0)$) --
    node[font=\scriptsize, text=black!60, fill=white, inner sep=1pt, left=2pt]
    {1. request ticket}
    ($(kdc.south)+(-1.20,0)$);
  \draw[ado arrow, draw=black!60, line width=1.05pt]
    ($(kdc.south)+(1.20,0)$) --
    node[font=\scriptsize, text=black!60, fill=white, inner sep=1pt, right=2pt]
    {2. issue service ticket}
    ($(requestor.north)+(1.20,0)$);

  \node[ado group, minimum width=14.20cm, minimum height=3.10cm, anchor=north] (services) at (0,-3.45) {};
  \node[ado label, font=\bfseries\normalsize] at ($(services.north)+(0,-0.38)$)
    {Protected Hadoop service layer};

  \node[ado box={5.25cm}{adobluefill}, minimum height=1.20cm, inner sep=5pt] (namenode) at (-3.25,-4.90)
    {\textbf{\small NameNode --- VM2}\\[2pt]
     \footnotesize service principal and keytab};
  \node[ado box={5.25cm}{adobluefill}, minimum height=1.20cm, inner sep=5pt] (datanodes) at (3.25,-4.90)
    {\textbf{\small DataNodes --- VM1/VM2/VM3}\\[2pt]
     \footnotesize service principals and keytabs};

  \node[ado annotation, font=\scriptsize, text=black!60] at (0,-6.08)
    {validated tickets enable client-to-service and service-to-service authentication};

  \draw[ado arrow, draw=black!60, line width=1.05pt]
    (requestor.south) --
    node[font=\scriptsize, text=black!60, fill=white, inner sep=1pt, right=2pt]
    {3. present ticket}
    (services.north);
\end{tikzpicture}

    \caption{Kerberos ticket flow for authenticated Hadoop clients and services.}
    \label{fig:chapter7-kerberos-identity-flow}
\end{figure}

\section{SPNEGO-Protected HDFS Browser Access}
The HDFS browser path is protected through SPNEGO layered on top of the Kerberos-based identity model. This means that access to the secured NameNode and DataNode web interfaces is not treated as a public or weakly protected administrative convenience. It becomes part of the authenticated architecture \cite{apache_hadoop_docs,mit_kerberos_docs}.

This protection extends the Hadoop security posture to the human-facing HDFS administration and browsing path instead of securing only back-end service interactions.

The implementation required both server-side and client-side alignment:
\begin{itemize}
    \item server-side Kerberos-aware Hadoop configuration;
    \item SPNEGO-protected web endpoints;
    \item HTTPS exposure of the HDFS browser interfaces;
    \item MIT Kerberos client configuration and a valid ticket on the Windows workstation;
    \item Firefox SPNEGO configuration aligned with the documented project hostname aliases rather than raw IP navigation.
\end{itemize}

This browser workflow is integration-sensitive because successful negotiation depends on the client ticket cache, hostname resolution, and Firefox configuration. A client-side SPNEGO failure must therefore be distinguished from the health of the Kerberos-secured Hadoop services themselves.

\begin{figure}[H]
    \centering
    \begin{tikzpicture}[node distance=0.6cm]
    \node[ado box=adograyfill, text width=10.4cm] (client)
    {\adoboxcontent{Windows client}{Kerberos ticket + Firefox SPNEGO configuration}};
    \node[ado box=adotealfill, text width=10.4cm, below=of client] (https)
    {\adoboxcontent{HTTPS HDFS browser endpoint}{SPNEGO-protected NameNode / DataNode web path}};
    \node[ado box=adobluefill, text width=10.4cm, below=of https] (hdfs)
    {\adoboxcontent{Secured Hadoop web services}{authenticated browsing of distributed file-system interfaces}};
    \draw[ado arrow] (client) -- (https);
    \draw[ado arrow] (https) -- (hdfs);
\end{tikzpicture}

    \caption{SPNEGO-authenticated HTTPS access to HDFS browser interfaces.}
    \label{fig:chapter7-hdfs-browser-spnego}
\end{figure}

\section{TLS-Protected Interfaces}
TLS is used to protect the main browser-facing interfaces of the platform. The objective is not to wrap every internal port indiscriminately, but to secure the access paths used for administration, monitoring, and analytical interaction. Transport protection complements identity protection rather than replacing it \cite{apache_hadoop_docs,trino_docs}.

Two main exposure zones are protected:
\begin{itemize}
    \item the HDFS browser interfaces through HTTPS;
    \item the Trino analytical access gateway through HTTPS.
\end{itemize}

The local control-plane interfaces are also exposed through a local TLS gateway, establishing a consistent protected entry path for Airflow, Grafana, Prometheus, and Kafka UI.

The deployment uses internal or self-signed trust paths rather than a public enterprise certificate chain. Browser warnings may therefore appear until the local certificate authority is trusted on the client machine. This limitation affects certificate trust management, not the presence of encrypted transport.

\begin{figure}[H]
    \centering
    \begin{tikzpicture}[node distance=0.7cm and 0.7cm]
    \node[ado box=adotealfill, text width=4.7cm] (hdfs)
    {\adoboxcontent{HDFS secure zone}{HTTPS browser access for NameNode / DataNode interfaces}};
    \node[ado box=adopurplefill, text width=4.7cm, right=of hdfs] (trino)
    {\adoboxcontent{Trino secure zone}{authenticated gateway over protected analytical access}};
    \node[ado box=adograyfill, text width=10.4cm, below=0.9cm of $(hdfs)!0.5!(trino)$] (cp)
    {\adoboxcontent{Local control-plane secure zone}{TLS-protected access to Airflow, Grafana, Prometheus, and Kafka UI}};
    \draw[ado dashed arrow] (hdfs) -- (cp);
    \draw[ado dashed arrow] (trino) -- (cp);
\end{tikzpicture}

    \caption{TLS-protected exposure zones for data-plane and control-plane interfaces.}
    \label{fig:chapter7-tls-exposure-zones}
\end{figure}

\section{Authenticated Trino Access}
Trino is not exposed as an open analytical endpoint. Instead, the project implements an authenticated gateway in front of the coordinator-side browser access path. This has two important effects \cite{trino_docs}:
\begin{itemize}
    \item it reduces uncontrolled direct access to the analytical query service;
    \item it applies an explicit authentication boundary to user-facing SQL access.
\end{itemize}

The gateway is hosted on VM2 together with the Trino coordinator, which is consistent with the node's role as the entry point for coordination-heavy services. Through this design, the browser-facing access path is separated from the internal worker topology. Users connect to the protected entry point, while the internal distributed Trino execution still takes place through the coordinator and worker nodes.

Authentication is enforced at the gateway level; this control should not be interpreted as a complete enterprise IAM or fine-grained authorization framework.

\begin{figure}[H]
    \centering
    \begin{tikzpicture}[node distance=0.65cm]
    \node[ado box={9.8cm}{adograyfill}] (user)
    {\adoboxcontent{Analytical users}{browser login through protected entry point}};
    \node[ado box={9.8cm}{adotealfill}, below=of user] (gateway)
    {\adoboxcontent{Authenticated gateway}{credential check + controlled external access path}};
    \node[ado box={9.8cm}{adopurplefill}, below=of gateway] (coord)
    {\adoboxcontent{Trino coordinator on VM2}{query planning + worker coordination}};
    \node[ado box={9.8cm}{adoorangefill}, below=of coord] (workers)
    {\adoboxcontent{Distributed workers}{query fragments executed on VM1 and VM3}};
    \draw[ado arrow] (user) -- (gateway);
    \draw[ado arrow] (gateway) -- (coord);
    \draw[ado arrow] (coord) -- (workers);
\end{tikzpicture}

    \caption{Authenticated Trino path from the VM2 gateway to distributed workers.}
    \label{fig:chapter7-trino-auth-gateway}
\end{figure}

\section{SSH-Based Remote Spark Submission}
The orchestration model relies on a remote Spark submission path from the control plane to the Spark submit client on VM2. This avoids fragile local-driver networking assumptions and keeps the operational launch path aligned with the distributed cluster environment.

From a security perspective, this means the orchestration node must not submit jobs through uncontrolled or manually repeated remote actions. The platform therefore uses an SSH-based remote execution path with dedicated key material and documented access preparation \cite{openssh_docs}.

This design improves the platform in three ways:
\begin{itemize}
    \item it keeps the control plane separate from the heavy Spark runtime;
    \item it avoids exposing the orchestration flow through weaker manual execution habits;
    \item it provides a consistent operational bridge between local Airflow logic and the distributed cluster leader on VM2.
\end{itemize}

This SSH path is therefore the authenticated operational bridge through which orchestration reaches distributed compute.

\begin{figure}[H]
    \centering
    \begin{tikzpicture}[x=1cm,y=1cm,font=\small]
  \node[ado box={11.80cm}{adograyfill}, minimum height=1.15cm, inner sep=5pt] (control) at (0,0)
    {\textbf{\normalsize Local Airflow control plane}\\[-1pt]
     \footnotesize creates a controlled remote Spark execution request};

  \node[ado box={11.80cm}{adotealfill}, minimum height=1.20cm, inner sep=5pt] (ssh) at (0,-1.75)
    {\textbf{\normalsize SSH-secured operational bridge}\\[-1pt]
     \footnotesize dedicated key material and documented VM2 access preparation};

  \draw[{Latex[length=2.6mm,width=1.6mm]}-{Latex[length=2.6mm,width=1.6mm]},
        draw=black!60, line width=1.0pt]
    (control.south) --
    node[font=\scriptsize, text=black!60, fill=white, inner sep=1pt, right=2pt]
    {remote request and task status}
    (ssh.north);

  \draw[draw=black!35, dashed, line width=0.8pt] (-6.65,-2.70) -- (6.65,-2.70);
  \node[font=\scriptsize, text=black!55, fill=white, inner sep=2pt] at (-4.25,-2.70)
    {local workstation / Linux data plane};

  \node[ado box={11.80cm}{adobluefill}, minimum height=1.20cm, inner sep=5pt] (submit) at (0,-3.60)
    {\textbf{\normalsize Spark submit client --- VM2}\\[-1pt]
     \footnotesize receives the remote command and launches the cluster application};

  \draw[{Latex[length=2.6mm,width=1.6mm]}-{Latex[length=2.6mm,width=1.6mm]},
        draw=black!60, line width=1.0pt]
    (ssh.south) --
    node[font=\scriptsize, text=black!60, fill=white, inner sep=1pt, right=2pt]
    {authenticated command / response}
    (submit.north);

  \node[ado group, minimum width=14.20cm, minimum height=3.85cm, anchor=north] (cluster) at (0,-4.65) {};
  \node[font=\bfseries\small, anchor=west] at ($(cluster.north west)+(0.35,-0.35)$)
    {Distributed Spark execution zone};

  \node[ado box={11.60cm}{adopurplefill}, minimum height=1.10cm, inner sep=5pt] (master) at (0,-5.75)
    {\textbf{\small Spark master --- VM2}\\[2pt]
     \footnotesize coordinates the application and schedules executor tasks};
  \node[ado box={11.60cm}{adoorangefill}, minimum height=1.10cm, inner sep=5pt] (workers) at (0,-7.40)
    {\textbf{\small Registered worker pool}\\[2pt]
     \footnotesize VM1, VM2, and VM3 capacity; active allocation depends on workload};

  \draw[{Latex[length=2.6mm,width=1.6mm]}-{Latex[length=2.6mm,width=1.6mm]},
        draw=black!60, line width=1.0pt]
    (submit.south) --
    node[font=\scriptsize, text=black!60, fill=white, inner sep=1pt, right=2pt]
    {cluster submission and exit status}
    (master.north);
  \draw[{Latex[length=2.6mm,width=1.6mm]}-{Latex[length=2.6mm,width=1.6mm]},
        draw=black!55, line width=0.95pt]
    (master.south) --
    node[font=\scriptsize, text=black!60, fill=white, inner sep=1pt, right=2pt]
    {tasks and execution status}
    (workers.north);
\end{tikzpicture}

    \caption{SSH-secured Spark submission from the local control plane to VM2.}
    \label{fig:chapter7-remote-spark-ssh}
\end{figure}

\subsection*{Operational Identity and Credential Separation}
The deployment separates non-interactive service identities from human-operated administration and demonstration identities. Kerberos principals and keytabs are assigned to protected Hadoop services, while browser credentials, gateway credentials, and SSH key material are scoped to their operational environments. This reduces uncontrolled credential sharing and clarifies which identity is used at each access boundary.

The implemented separation is a baseline control rather than a centralized secret-management system. Automated secret rotation, enterprise vault integration, and organization-wide identity federation remain outside the current scope.

\section{Encrypted Runtime Storage Support}
The project also implements optional support for encrypted runtime storage through LUKS on the Linux virtual machines. This measure targets \textbf{data at rest} rather than service authentication or traffic encryption \cite{cryptsetup_docs}.

When enabled, this control protects runtime directories used by Kafka, HDFS, Hive, Spark, and Trino on VM1, VM2, and VM3. It complements the controls applied to service identity, browser access, and network transport.

The implemented logic therefore complements the other controls:
\begin{itemize}
    \item Kerberos protects identity,
    \item SPNEGO protects browser negotiation,
    \item TLS protects exposed traffic,
    \item the authenticated Trino gateway protects analytical access,
    \item SSH protects remote Spark submission,
    \item LUKS-backed storage protects persisted runtime data when enabled.
\end{itemize}

This layered view helps the report explain that security was approached as a set of complementary boundaries rather than a single technology checkbox.

\begin{figure}[H]
    \centering
    \begin{tikzpicture}[x=1cm,y=1cm,font=\small]
  \node[ado group, minimum width=14.60cm, minimum height=6.35cm, anchor=north] (coverage) at (0,0) {};
  \node[ado label, font=\bfseries\large] at ($(coverage.north)+(0,-0.38)$)
    {Implemented security coverage};
  \node[ado annotation, font=\scriptsize, text=black!55] at (0,-0.84)
    {complementary controls protecting distinct risk surfaces};
  \node[ado annotation, font=\scriptsize, text=black!55] at (0,-1.18)
    {supported by distinct service, administration, and demonstration identities};

  \node[ado box={5.70cm}{adotealfill}, minimum height=1.05cm, inner sep=5pt] (identity) at (-3.35,-2.00)
    {\textbf{\small Identity authentication}\\[2pt]
     \footnotesize Kerberos identities for Hadoop};
  \node[ado box={5.70cm}{adobluefill}, minimum height=1.05cm, inner sep=5pt] (browser) at (3.35,-2.00)
    {\textbf{\small Browser authentication}\\[2pt]
     \footnotesize SPNEGO for HDFS web access};

  \node[ado box={5.70cm}{adopurplefill}, minimum height=1.05cm, inner sep=5pt] (transport) at (-3.35,-3.65)
    {\textbf{\small Transport encryption}\\[2pt]
     \footnotesize TLS for HDFS, Trino, and local UIs};
  \node[ado box={5.70cm}{adotealfill}, minimum height=1.05cm, inner sep=5pt] (analytical) at (3.35,-3.65)
    {\textbf{\small Analytical access}\\[2pt]
     \footnotesize authenticated Trino gateway};

  \node[ado box={5.70cm}{adograyfill}, minimum height=1.05cm, inner sep=5pt] (remote) at (-3.35,-5.30)
    {\textbf{\small Remote operations}\\[2pt]
     \footnotesize SSH remote Spark submission};
  \node[ado box={5.70cm}{adoorangefill}, minimum height=1.05cm, inner sep=5pt] (storage) at (3.35,-5.30)
    {\textbf{\small Data at rest}\\[2pt]
     \footnotesize optional LUKS-backed storage support};

\end{tikzpicture}

    \caption{Implemented security coverage across complementary protection dimensions.}
    \label{fig:chapter7-security-coverage}
\end{figure}

\section{Security Boundaries and Remaining Limits}
The implemented security layer provides a meaningful distributed-platform hardening baseline; it does not constitute a finished enterprise security program.

The current implementation can validly claim:
\begin{itemize}
    \item Kerberos-secured Hadoop runtime;
    \item SPNEGO-protected HDFS browser access;
    \item TLS-protected browser-facing endpoints;
    \item authenticated Trino access;
    \item SSH-based remote Spark submission;
    \item distinct service, administration, and demonstration identities;
    \item environment-scoped credential handling;
    \item optional encrypted runtime storage support on the Linux nodes.
\end{itemize}

At the same time, some limits remain outside the strict PFE scope:
\begin{itemize}
    \item enterprise-scale PKI and certificate lifecycle automation;
    \item full centralized secret management;
    \item fine-grained authorization frameworks across all analytical consumers;
    \item large-scale security auditing, SIEM integration, or intrusion detection;
    \item organization-wide identity federation.
\end{itemize}

These boundaries distinguish the controls implemented in the project from future enterprise hardening work.

\begin{figure}[H]
    \centering
    \begin{tikzpicture}[x=1cm,y=1cm,font=\small]
  \node[ado group, minimum width=14.60cm, minimum height=6.10cm, anchor=north] (scope) at (0,0) {};
  \node[ado label, font=\bfseries\large] at ($(scope.north)+(0,-0.38)$)
    {Security scope boundary};

  \node[ado box={5.85cm}{adotealfill}, minimum height=4.55cm, inner sep=7pt] (implemented) at (-3.55,-3.35)
    {\textbf{\normalsize Implemented baseline}\\[-1pt]
     \scriptsize completed project scope\\[6pt]
     {\scriptsize
     \renewcommand{\arraystretch}{1.18}
     \begin{tabular}{@{}l@{}}
       \(\bullet\) Kerberos principals and keytabs \\
       \(\bullet\) SPNEGO + HTTPS for HDFS \\
       \(\bullet\) TLS-protected browser gateways \\
       \(\bullet\) Authenticated Trino gateway \\
       \(\bullet\) SSH remote Spark submission \\
       \(\bullet\) Role-separated identities \\
       \(\bullet\) Optional LUKS storage support
     \end{tabular}}};

  \node[ado box={5.85cm}{adograyfill}, minimum height=4.55cm, inner sep=7pt] (future) at (3.55,-3.35)
    {\textbf{\normalsize Future enterprise hardening}\\[-1pt]
     \scriptsize outside the PFE scope\\[6pt]
     {\footnotesize
     \renewcommand{\arraystretch}{1.24}
     \begin{tabular}{@{}l@{}}
       \(\bullet\) PKI and certificate automation \\
       \(\bullet\) Centralized secret management \\
       \(\bullet\) Fine-grained authorization \\
       \(\bullet\) Organization-wide federation \\
       \(\bullet\) SIEM, auditing, and detection
     \end{tabular}}};

  \draw[draw=black!35, dashed, line width=0.8pt] (0,-1.05) -- (0,-5.65);
  \node[font=\scriptsize, text=black!55, fill=white, inner sep=2pt, rotate=90] at (0,-3.35)
    {scope boundary};
\end{tikzpicture}

    \caption{Implemented security baseline and enterprise controls outside the PFE scope.}
    \label{fig:chapter7-security-boundaries}
\end{figure}

\section{Conclusion}
This chapter demonstrated a layered security baseline comprising Kerberos principals and keytabs, SPNEGO/HTTPS access to HDFS, TLS-protected browser gateways, authenticated Trino entry, SSH-based Spark submission, separated operational identities, and optional LUKS-backed storage. These controls harden the distributed deployment without being presented as a complete enterprise IAM program. Chapter VIII examines the resulting readiness, pipeline, observability, and security evidence.
